\section{Estrategia empírica}\label{sec:strategy}

\subsection{Resultados potenciales y el estimando}

Sea $i$ el índice de individuos, $g\in\{\Low,\mathrm{High}\}$ su grupo de calificación y
$t$ el trimestre. La reforma entra en vigor en el segundo trimestre de 2022, que denoto
$t^{\ast}$;\footnote{Como la reforma rigió desde el 1 de mayo, el segundo trimestre de 2022
contiene un mes previo y dos meses posteriores a la reforma, y los ingresos declarados se
refieren al mes anterior a la entrevista, de modo que incluso las entrevistas de mayo
reportan ingresos previos a la reforma (de abril). Por ello excluyo 2022Q2 de la
especificación base como trimestre de transición parcialmente tratado y lo reincorporo en
una prueba de robustez; el primer trimestre posterior a la reforma en la especificación
base es 2022Q3.} defino el indicador posterior a la reforma
$\Post_t=\mathbf{1}\{t\ge t^{\ast}\}$ y el indicador de tratamiento
$\Low_i=\mathbf{1}\{g_i=\Low\}$. Siguiendo el marco de resultados potenciales, sean
$Y_{it}(1)$ e $Y_{it}(0)$ el resultado del individuo $i$ en el trimestre $t$ con y sin la
reforma vigente. El parámetro de interés es el efecto promedio de la reforma sobre los
trabajadores poco calificados, respecto del contrafactual en el que habrían experimentado
la misma evolución que los calificados,
\begin{equation}
\mathrm{ATT} \;=\; \mathbb{E}\!\left[\,Y_{it}(1)-Y_{it}(0)\;\middle|\;\Low_i=1,\;t\ge t^{\ast}\right].
\end{equation}
Los datos son cortes transversales repetidos y no un panel, de modo que este objeto se
identifica a partir de comparaciones de medias de grupo en el tiempo y no de cambios dentro
de cada persona.

Como el salario mínimo del Perú es nacional y cambia en una sola fecha, no hay variación en
el \emph{momento} del tratamiento: todos los trabajadores quedan expuestos a la misma
política en el mismo instante. La identificación descansa, en cambio, en la variación de la
\emph{intensidad} de la exposición entre grupos de calificación. El diseño es, por tanto, una
diferencias en diferencias canónica de dos grupos con muchos periodos, y no un diseño de
adopción escalonada. Esta distinción importa para la elección del estimador. La literatura
econométrica reciente ha mostrado que el estimador de efectos fijos bidireccionales puede
estar severamente sesgado cuando el momento del tratamiento es escalonado y los efectos son
heterogéneos \citep{goodmanbacon_2021, dechaisemartin_2020, sun_abraham_2021,
callaway_santanna_2021}. Con una fecha de tratamiento común, sin embargo, estos problemas no
se presentan: el estimador de efectos fijos bidireccionales y su análogo de estudio de
eventos recuperan un promedio convexo y correctamente ponderado de los efectos del
tratamiento \citep{roth_2023_trending, baker_larcker_wang_2022}. Aun así, presento el
estimador doblemente robusto de \citet{santanna_zhao_2020} como complemento.

\subsection{Ecuaciones de estimación}

Mi especificación base es el modelo de diferencias en diferencias con efectos fijos
bidireccionales
\begin{equation}
Y_{idt} \;=\; \beta\,(\Low_i\times\Post_t) \;+\; \gamma\,\Low_i \;+\; \lambda_t \;+\; \mu_d
\;+\; X_{it}'\theta \;+\; \varepsilon_{idt},
\label{eq:twfe}
\end{equation}
donde $d$ denota el departamento, $\lambda_t$ son efectos fijos de trimestre, $\mu_d$ son
efectos fijos de departamento y $X_{it}$ es un vector de controles que incluye la edad, su
cuadrado, el sexo, un indicador de zona urbana, el estado civil y los años de educación (que
recogen la variación del nivel educativo dentro de cada grupo de calificación). El
coeficiente $\beta$ es la estimación de diferencias en diferencias. Los efectos fijos de
trimestre absorben cualquier choque común a ambos grupos en un trimestre dado, incluida la
inflación agregada, de modo que las estimaciones en logaritmos no dependen de si los
salarios se expresan en términos nominales o reales deflactados a nivel nacional. Las
observaciones se ponderan con los factores de expansión, y los errores estándar se agrupan
por departamento.

Para trazar la dinámica del efecto y contrastar el supuesto de identificación, estimo la
especificación dinámica de estudio de eventos
\begin{equation}
Y_{idt} \;=\; \sum_{k\neq -1}\beta_k\,\bigl(\Low_i\times\mathbf{1}\{t-t^{\ast}=k\}\bigr)
\;+\; \gamma\,\Low_i \;+\; \lambda_t \;+\; \mu_d \;+\; X_{it}'\theta \;+\; \varepsilon_{idt},
\label{eq:event}
\end{equation}
donde $k$ indexa los trimestres relativos a la reforma y el trimestre inmediatamente anterior
a ella ($k=-1$, el primer trimestre de 2022) es la categoría de referencia omitida. Los
coeficientes $\beta_k$ para $k<0$ miden tendencias previas diferenciales entre los grupos de
calificación y ofrecen una prueba directa de tendencias paralelas; los coeficientes para
$k\ge 0$ trazan la trayectoria del efecto.

\subsection{Supuestos de identificación y amenazas}

La identificación del ATT a partir de la ecuación~\eqref{eq:twfe} requiere que, en ausencia
de la reforma, los resultados de los trabajadores poco calificados y calificados hubieran
seguido trayectorias paralelas, condicionales a los controles y efectos fijos. Tres amenazas
merecen atención.

Primero, las \emph{tendencias diferenciales}. La educación está correlacionada con el sector,
la región y la formalidad, y los dos grupos de calificación podrían haber seguido
trayectorias divergentes por razones ajenas al salario mínimo. Evalúo esto directamente con
los coeficientes del estudio de eventos previos a la reforma y una prueba conjunta de que son
iguales a cero, y examino la robustez a tendencias lineales específicas por departamento.
Como mi ventana se abre en plena recuperación de la pandemia, cuando el empleo de los
trabajadores poco calificados y calificados se recuperó a ritmos distintos, presto especial
atención al primer trimestre de 2021 y presento especificaciones que lo excluyen.

Segundo, la \emph{sensibilidad a violaciones residuales}. Aun cuando las tendencias previas
son estadísticamente planas, no tienen por qué ser exactamente cero. Por ello presento el
análisis de sensibilidad de \citet{rambachan_roth_2023}, que pregunta cuán grande tendría que
ser una violación de tendencias paralelas posterior a la reforma, en relación con la mayor
violación observada antes de ella, para revertir cada conclusión.

Tercero, los \emph{efectos de derrame y de equilibrio general}. Si la reforma desplaza a
trabajadores poco calificados hacia el sector informal, podría deprimir indirectamente los
salarios en ese sector, y si los trabajadores calificados son sustitutos imperfectos, sus
resultados también podrían moverse y contaminar el grupo de comparación. No puedo descartar
esta posibilidad, e interpreto las estimaciones como efectos relativos. El grupo de
comparación no está del todo libre de exposición: el 40 por ciento de los asalariados
privados calificados percibía menos que el nuevo piso antes de la reforma, frente al 57 por
ciento de los poco calificados. En primer orden, esta exposición imperfecta atenúa las
estimaciones hacia cero (el diseño identifica el efecto de la \emph{brecha} de exposición de
diecisiete puntos, por lo que subestima el efecto sobre los plenamente expuestos), pero
reconozco que también convierte las estimaciones en objetos relativos y no absolutos. Por esta
razón complemento el contraste por educación con un diseño de exposición regional continua
que no depende en absoluto del grupo de comparación (Sección~\ref{sec:robustness}).

\subsection{Inferencia}

El tratamiento se asigna a un nivel agregado, el grupo de calificación, lo que plantea la
preocupación habitual de que los errores estándar convencionales sobrestimen la precisión
\citep{bertrand_2004}. La abordo de varias formas complementarias. Mi especificación base
agrupa los errores por departamento,\footnote{Los 25 conglomerados son los 24 departamentos
del Perú más la Provincia Constitucional del Callao. Con pocos conglomerados, los errores
estándar robustos por conglomerado convencionales pueden subestimar la incertidumbre, lo que
motiva las correcciones adicionales que se describen a continuación.} lo que permite una
correlación arbitraria dentro de cada departamento entre individuos y trimestres; presento
una alternativa que agrupa por departamento y grupo de calificación; y aplico la corrección
de muestra pequeña CR2 de reducción de sesgo de \citet{pustejovsky_tipton_2018} a celdas
agregadas a nivel departamental, donde resulta computacionalmente factible. Como estas
correcciones pueden seguir siendo optimistas en un diseño de dos grupos, las complemento con
inferencia por aleatorización a nivel departamental, permutando la incidencia regional del
salario mínimo entre departamentos para construir una distribución nula exacta para el
diseño de intensidad de exposición.
